\section{Causal certificate controls on an unchanged path}
\label{sec:tangentmonitors}
We describe two standard lower-bound constructions used as stopping
controls. Neither is a new primal--dual method. Write
$g_j=\nabla f(x_j)$ and $a_j=f(x_j)-\inner{g_j}{x_j}$.
For a completed block $B$ of $K$ consecutive source observations, set
\[
 \bar a_B=K^{-1}\sum_{j\in B}a_j,\qquad
 \bar g_B=K^{-1}\sum_{j\in B}g_j,\qquad
 b_B=\bar a_B+\min_{v\in Q}\inner{\bar g_B}{v}.
\]
The last minimization costs one additional full LMO. Disjoint blocks
start at source $j=0$; an incomplete block is not queried.

\begin{proposition}[Elementary stopping control]\label{prop:blocks}
Let $f(x_k)$ be nonincreasing along a fixed feasible path, and let $\ell_k$
be the accumulated scalar tangent lower bound. Define
\[
 \widetilde\ell_k=\max\{\ell_k,\ b_B:B\text{ completed by }k\},
 \qquad \widetilde U_k=f(x_k)-\widetilde\ell_k.
\]
Then $\delta_k\le\widetilde U_k\le U_k$, and
$\widetilde U_k$ is nonincreasing. Its first crossing of a prescribed
tolerance occurs no later on this path than that of $U_k$.
Up to source $k$, exactly $\lfloor(k+1)/K\rfloor$ extra LMOs are used,
with no extra function or gradient calls.
\end{proposition}
\begin{proof}
Convexity gives $f(v)\ge a_j+\inner{g_j}{v}$ for each source and every
$v\in Q$. Averaging and minimizing gives $b_B\le f_Q^*$.
Taking the maximum preserves validity and is nondecreasing. Subtracting
it from the nonincreasing function values proves the assertions;
the oracle count follows from the number of completed blocks.
\end{proof}

In particular, if the original algorithm stops after $N$ updates, this
control uses at most $N+1+\lfloor(N+1)/K\rfloor$ total LMOs and no more
function/gradient calls than the original run. This is a pathwise
comparison, not an improved worst-case convergence rate. Earlier stopping
need not reduce the total LMO count. The monitor affects neither the
line search nor the initial curvature estimate.

Average tangent bounds need not dominate the best scalar tangent bound.
For $f(x)=x^2/2$ on $[-1,1]$, the two-source histories $(-1,1)$ and
$(-1,-1/2)$ give block bounds $-1/2$ and $-17/16$, versus scalar maxima
$-3/2$ and $-5/8$, respectively. The maximum in
Proposition~\ref{prop:blocks} prevents deterioration. For a globally
defined differentiable convex $f$, Fenchel equality and Jensen give
\[
 b_B=-K^{-1}\sum_{j\in B}f^*(g_j)-\sigma_Q(-\bar g_B)
 \ \le\ -f^*(\bar g_B)-\sigma_Q(-\bar g_B)=\psi(\bar g_B).
\]
Here $f^*$ is the convex conjugate, distinct from the scalar optimum
$f_Q^*$. Pedregosa et al.~\cite[Theorem 3]{pedregosa} already analyze an
averaged-gradient Fenchel dual for adaptive FW variants, with their
specified weights. Our uniform, disjoint-block tangent computation
does not evaluate that dual or inherit its theorem.

A separate control receives a valid strong-convexity constant $\mu>0$.
Minimizing the quadratic lower model over the ambient space gives
\[
 b_j^\mu=f(x_j)-\frac{\norm{g_j}^2}{2\mu}\ \le f_Q^*,\qquad
 \widetilde\ell_k^\mu=\max\{\ell_k,\max_{j\le k}b_j^\mu\}.
\]
This standard bound is recorded, for example, in
\cite[Appendix B]{karimi}; it remains valid for a constrained minimum
because minimizing the model over $\mathbb R^d$ only lowers it.
It requires no extra LMO but uses more information than the block
controls. At a boundary solution a nonzero gradient can keep it weak.
All assertions here use exact arithmetic; the experiment audits
floating-point bounds against independent reference intervals.
